\section{Existence of maps in the original theta datum}

% BEGIN E77A1
\paragraph{E77A1. The projection at the bottom of the tower.}
Fix one original real datum $R$ with its level, family, models and
group record. Assume its weak chain and weak row theta conditions.
These hypotheses hold for the real part of a datum with the source's
weak characterization at the specified face. For a twist $t$, compose
the maps already stored in the chain from $H_X(t)$ down to its zeroth
multiplicity space. Convert the tensor-linear output to a multilinear
map and denote it by $D_t(x)$. This is legitimate even before any
integral has been identified with a chain map. Each arrow has the
next space as its codomain, so induction on the number of arrows puts
the composite in that zeroth multiplicity space. The rank of $Y_0$
is precisely $r$. At every extra twist prime the second diagonal
block has size zero and determinant one, and the flag with one block
of size $r$ imposes no additional condition. Thus both the level and
its character there agree with the untwisted ones. In particular
$D_t(x)$ lies in the principal flag multiplicity space $M$ at the
trivial character in the base datum.

Let $c_a$ be the coefficients of the recorded good-series idempotent
transported to the finite group $G$ at the base embedding. Write $T_a$
for its flag action on $M$. To see that these are actions on the stated
space, lift the prescribed leading block at each marked split place.
In the lattice frame lift a residue matrix in $\mathrm{GL}_r$ to a
local integral invertible matrix (a lift of its determinant is a local
unit). Use its block diagonal extension by the identity, put the
inverse Hermitian adjoint at the conjugate place, and put the identity
away from the pair. Self-duality of the completed lattice makes the
conjugate component a lattice stabilizer as well. The chosen primes
and their conjugates are disjoint. Indeed, conjugation fixes the
underlying rational prime, and above each of the two distinct split
rational primes the marked divisors use only one side. The local
matrices therefore multiply to a simultaneous lift. Two lifts differ
by the principal flag subgroup. Its character cancels the ambiguity
in the character-normalized right translation, giving $T_aT_b=T_{ab}$
on this type. A finite right translation preserves rational invariance,
cuspidality, continuity and archimedean equivariance. The principal
flag level and its character are preserved by the normalization.
The same statements pass to the linear span defining $M$.

Set $P=\sum_a c_a T_a$. The coefficient of $d$ in the square of the
group algebra idempotent is $\sum_a c_a c_{a^{-1}d}=c_d$. Hence
\[
 P^2D_t(x)=\sum_d\Big(\sum_a c_a c_{a^{-1}d}\Big)T_d D_t(x)
          =P D_t(x).
\]
Here right translation on a multilinear map acts on its values.
The good multiplicity space is $M\cap\ker(P-1)$. Its Hilbert model
has a linear injection $e_E$ onto this precise subspace. Thus
$\Theta_t(x)=(e_E|^{\mathrm{im}(e_E)})^{-1}P D_t(x)$ exists and
is linear. This gives exactly $R.\mathrm{ex\_Th}$, including its
quantifier over all twists.
% END E77A1

% BEGIN E77A2
\paragraph{E77A2. The projected row map from its unprojected companion.}
Fix a row $i$, the distinguished embedding $\sigma_0$ and a twist $t$.
Let $e_W:H_W(i,\sigma_0)\to M_W$ be the stored realization.
The model equations give $\mathrm{im}(e_W)=W_{\rm all}$, injectivity,
and $e_W a_g=T_g e_W$ for the stored linear actions $a_g$. Assume
that the unprojected stored map $F=\mathrm{thF}_{i,t}$ takes values
in $W_{\rm all}$. Define
\[
 G_{i,t}(x)=\sum_{g\in G}c_g a_g\,
                  (e_W|^{\mathrm{im}(e_W)})^{-1}F(x).
\]
Applying $e_W$ term by term gives $\sum_g c_g T_g F(x)$, exactly
postcomposition with the specified row projector. If another map
satisfies that identity, injectivity gives equality on every $x$.
This constructs the native $\mathrm{PThG}$ without a choice of
preimages for individual vectors. For all $i,t$, the image assumption
follows from $R.\mathrm{ex\_thF}$ and the weak row equations:
the existence of a witness makes the stored $F$ satisfy the entire
$\mathrm{PThF}$ predicate, whose first conjunct is this membership.
Thus $R.\mathrm{ex\_thF}\Rightarrow R.\mathrm{ex\_thG}$ under the
stated model and weak conditions.
% END E77A2

% BEGIN E77A3
\paragraph{E77A3. The two maps at the definite end of the tower.}
Put $n=r+(m-1)t_0$ and $N=n+t_0$, where $t_0$ denotes the fixed
rank increment. Both terminal spaces have positive signature at
every real place. Their smooth representations are finite-dimensional
irreducible compact types. In the source's decreasing-partition
convention their weights, after the indicated determinant twists,
complex conjugations and reversals, are a partition $\nu$ for $U(N)$
and its restriction to $n$ entries for $U(n)$, with $\nu_j=0$ for
$j\ge n$ (numbering entries from zero). This follows by specializing
the chain parameters at $k+1=m$, so that both negative index blocks
are empty. In the source notation their actual entries are
$\Lambda_a=x_3-\nu_{n-1-a}$ and
$\Lambda'_b=-x_1-\nu_{N-1-b}$. At each place use the polynomial decomposition
\[
 \operatorname{Sym}(\mathbb C^N\otimes\mathbb C^n)
   =\bigoplus_{\ell(\lambda)\le n}
       S_\lambda(\mathbb C^N)\otimes S_\lambda(\mathbb C^n).
\]
Here the standard actions are on the polynomial variables; in a
coordinate-function convention the corresponding factors are duals.
The equality follows degree by degree by applying Schur--Weyl duality
to both tensor powers and taking diagonal symmetric-group invariants.
The equal partitions pair with multiplicity one. In particular the
$\nu$ summand supplies a nonzero polynomial map
$\Psi:E\otimes H_n\to\mathcal P$ for the transformed types, where
$J:H_N\to E$ is the conjugate-linear identification specified by the
twists. With the Fischer form linear in its first slot define
\[
 \langle u,I(P,v)\rangle=[P,\Psi(Jv,u)]_{\rm F}.
\]
The two conjugations on the right make $I$ linear in $v$. Positivity
of the Fischer form on a nonzero $\Psi(Jv,u)$ gives $I\ne0$. Unitary
row and column substitutions preserve the form (in the monomial
basis the squared norm of $z^a$ is $a!$). Moving them across the form
gives the column covariance and the simultaneous row invariance,
including the determinant twists already built into $\Psi,J$.
Tensor the $I$ over the real places, identifying a product of
polynomials in disjoint variables with the tensor of its factors.
A tensor of nonzero local values is nonzero. For each fixed polynomial
the resulting map in the vector variable is bounded in these finite
dimensions. Unitary substitution preserves the finite-dimensional
degree filtration, so differentiation of the two group identities
gives the Lie identities by the product rule. The realization on the
nonzero Fock vacuum is injective on each degree space and intertwines
these substitutions for every weak terminal vacuum and its Weil
datum. In particular, a prescribed pointwise derivative that is again
a realized polynomial agrees with the derivative in the degree space:
finitely many point evaluations form a separating family on that
finite-dimensional image. The same $I$ works for all twists. This
supplies $R.\mathrm{ex\_ItX}$. The weak condition now certifies
the operator actually stored in $R$.

For this stored operator write $I$ again, and let $C$ be the range
of the source coefficient map $c_X.\mathrm{toMap}(\mathrm{pr}_X)$
composed with the realization $e_X$. For $F\in C$ put
\[
 B(P,v)(F)(g)=\int_{\mathcal F_X}
      \theta(g,h;\mathrm{pt}_X(P))\,F(v)(h)\,d\nu_X(h).
\]
The definite quotient is compact, the kernel and the coefficient form
are continuous, and both descend in the integration variable as
required. The product is integrable on the chosen fundamental domain,
and the result is continuous in $g$ by compact domination locally in
$g$. Haar invariance gives simultaneous row invariance of $B$. The
nonzero determinant scalars in the substitution law can be absorbed
into the type actions.

Here is why the integral factors through the \emph{stored} $I$.
After linearization in $P\otimes v$, the dual image of $I$ is a
nonzero column-stable subspace of the row-invariant scalar
functionals. That last space is irreducible for the column action:
in the displayed decomposition the row type singles out exactly one
partition, paired with a single irreducible column type, and the same
argument takes tensor products over the places. Scalar invariant forms
on the full polynomial algebra are supported in the degree singled out
at each place by the central circle of the row type, so algebraic duals
introduce no additional summands here. Every functional $\ell B$
therefore equals $\ell' I$. Linear functionals separate vectors,
whence $\ker I\subseteq\ker B$. Also $\mathrm{im}(I)$ is nonzero and
column-stable and so is all of the irreducible target. Thus $B=L I$
for a linear $L$. On generators $I(P,v)$, the finite level character
follows by moving the finite translation to the vacuum, rational
invariance follows from that of the kernel, and archimedean covariance
follows from column covariance. These statements and continuity
extend by linearity. The definite target has no proper rational
parabolic, so cuspidality adds no further condition. Evaluation
$F\mapsto (y\mapsto L(y)(F))$ consequently lands in the required
terminal multiplicity space, with exactly the displayed integral and
its integrability. This proves $R.\mathrm{ex\_ioX}$ for every twist.
% END E77A3

% BEGIN E77A4
\paragraph{E77A4. One joint lowest-type map for every twist of a row.}
Write $p_\tau+q_\tau=\mathrm{rank}(W_i)\le\mathrm{rank}(V_i)$
and let $\Delta_i$ be the prescribed product of flag minors. The
polynomial space carries commuting actions of the coefficient-side
unitary group $U=\prod_\tau U(\mathrm{rank}(V_i))$ and the target
compact group $K=\prod_\tau U(p_\tau)\times U(q_\tau)$, with the
determinant characters of the row. Reverse the order of the negative
rows using the permutation $\kappa\in K$ and let $\Delta'$ be its
twisted action on $\Delta_i$. It is a simultaneous bottom vector for
the two polynomial actions in the source's holomorphic conventions.
Indeed each triangular substitution multiplies a flag minor by the
product of the diagonal entries on its flag. Taking differences of
successive partition entries as minor exponents makes the powers
telescope. For the $U$ action the resulting weight is $-\lambda$,
where $\lambda$ is the weight of the coefficient model. For the
holomorphic $K$ action it is $-\mathrm{lowK}$; this latter action uses
inverse transpose on the compact matrix. These are the original
row-weight equalities, with determinant characters included.
Each polynomial factor is nonzero and the polynomial ring is a domain.

Realize the complex coefficient space as the frame-polynomial
module $V_\lambda$. An extremal vector $w_0$ generates it under $U$.
There is a $U$-map taking $w_0$ to $\Delta'$ in the polynomial space.
For example, decompose the polynomial space of degree at most
$\deg\Delta_i$ orthogonally under the compact group. Each nonzero
projection of $\Delta'$ has the same bottom weight and hence belongs
to a copy of the same irreducible representation. In each copy take
the equivariant map that sends $w_0$ to that component and sum the
maps. All values $\psi(w)$ are in the same $K$-bottom space because
the two actions commute. On this bottom space extend
$x_0\otimes y\mapsto y$ equivariantly in the lowest $K$-type
variable. More explicitly, on each isotypic summand $L\otimes M_L$
with bottom line $\mathbb C b\subset L$, identify $b\otimes u$ with
$u$ and send $(x,b\otimes u)$ to $x\otimes u$ in that summand.
Use an equivariant projection to a chosen copy of the type in the
source if the isotypic space has multiplicity. Every operator commuting
with $K$ acts on the multiplicity factor, so this extension also
commutes with $U$. Pull back the $K$ action by the compact
conjugate/inverse-transpose convention to get the physical substitution
law in the actual frame. Denote the resulting bilinear map by $j_i$.

Choose a pure generator $\otimes_\tau v_\tau$ whose compact translate
acts by $\kappa$ on $\Delta'$; the factors can be chosen separately at
the real places. The square of the twisted action of $\kappa$ is the
identity. The value of $j_i$ at this generator and $w_0$ is therefore
$\Delta_i$ itself. Rescale the inputs to unit norm and multiply $j_i$
by $\|w_0\|\prod_\tau\|v_\tau\|$ to keep that value fixed. The two
compact covariance laws for the polynomial actions pass to each
realization $\mathrm{pt}_{i,t}$ by its vacuum substitution formula.
The weights, the two compact models and the construction of $j_i$
preceded the choice of $t$. This is the order of quantifiers in
$R.\mathrm{ex\_ioP}$, including the exact unit-vector normalization.
% END E77A4

% BEGIN E77A5
\paragraph{E77A5. Extending a theta integral from the lowest type.}
For the row of E77A4 the weak characterization makes the stored
$\mathrm{ioP}_i$ satisfy both covariance laws and attain $\Delta_i$.
Fix a twist and let $F_x$ be the coefficient-valued form associated
to $x\in H_V(i,t)$. For a vector $v$ in the tensor of lowest types
put
\[
 U_x(v)(g)=\int_{[U(V_i)]}\theta\big(h,g;
       \mathrm{pt}_{i,t}(\mathrm{ioP}_i(v,F_x(h)))\big)\,d\nu(h).
\]
The quotient has the measure and fundamental domain of $R$. Its
definiteness gives compactness. The integrand descends in $h$ by
the coefficient and theta transformation laws and the integral is
bilinear in $x,v$. Translation by the target principal level fixes the
finite vacuum, while the compact covariance of $\mathrm{ioP}_i$
gives precisely the lowest-type law on the right. Here are the
additional analytic facts about this particular integral that enter
the extension. The source proves smoothness and uniform moderate growth
on the target by its definite theta kernel estimates. At every rational
isotropic line it supplies a finite nonempty family of Heisenberg
translates whose sum on the row test vanishes. Averaging over the
corresponding unipotent quotient kills the constant term. On the
archimedean polynomial the mixed Laplacians annihilate $\Delta_i$:
the two systems of minors use disjoint end sets of columns, since
$p_\tau+q_\tau\le\mathrm{rank}(V_i)$. The oscillator lowering formula
turns this into annihilation by the antiholomorphic tangent directions.
Equivariance and generation by the compact types carry this annihilation
to the row integrals. These kernel estimates, the finite-place
vanishing, and the holomorphic discrete-series realization are inputs
at this step.

Use the rapid-decay theorem for smooth cuspidal forms of uniform
moderate growth. The function attached to each extremal pure
lowest-type vector is bounded and hence square integrable on the
finite-volume target quotient. Its level character, compact type
and lowering equations are those of the specified holomorphic
discrete series. The bounded holomorphic realization theorem applied
to this vector gives a multiplicity map with that value. Compact
equivariance propagates equality from this tuple to the tensor of
the lowest types: each local lowest type is irreducible here, its
nonzero extremal vector is cyclic, and multilinearity extends equality
from the product of the orbits to their span. Thus for each $x$ some
element of $W_{\rm all}$ restricts to $U_x$. To choose linearly in $x$,
restrict $W_{\rm all}$ to the space of multilinear functions on the
lowest types and take a linear section on the range. Composing this
section with $x\mapsto U_x$ gives $\mathrm{thF}'_{i,t}$. The companion
map on $\mathrm{im}(\mathrm{extV}_{i,t})$ is the same integral on a
coefficient-valued input form, so it is well defined even if two input
vectors yield the same form. This gives the range, integrability and
factorization clauses of $\mathrm{PThF}$. In particular the original
definite-row kernel and realization theorems, together with
\emph{CuspRapidDecay}, suffice for $R.\mathrm{ex\_thF}$ and hence for
$R.\mathrm{ex\_thG}$ by E77A2.
% END E77A5

% BEGIN E77A6
\paragraph{E77A6. What remains of the eight existence clauses.}
Under the weak, level, model and group-record hypotheses of E77A1,
and with the definite-row analytic inputs of E77A5, one has
\[
 R.\mathrm{ExAt}\quad\Longleftrightarrow\quad
      R.\mathrm{ex\_ItY}\ \land\ R.\mathrm{ex\_ioY}.
\]
Projection of the eight-field record proves the forward implication.
In the other direction insert the two indicated witnesses together
with E77A3 for its terminal entries, E77A1 for the bottom projection,
E77A4 for the row polynomial map and E77A5 and E77A2 for the two row
multiplicity maps. Every entry belongs to the same $R$.

For precision about the remaining quantifiers, at each $k+1<m$ the
intertwiner must be chosen \emph{before} the twist; the multiplicity
lift is required for each $(t,k)$, and uses the \emph{stored}
intertwiner. In the formal construction the model for the
archimedean covering discrete series and its Harish--Chandra module
supply a positive doubling pairing on joint harmonics and an embedding
of the discrete-series tensor in Fock space. Together with the vacuum
isometry and character these realize the first common intertwiner
on models with the prescribed parameters. The chain parameters follow
from the host's node gap. Existence permits the weak characterization
to certify the stored intertwiner. The inner multiplicity construction
also uses rapid decay for the cuspidal $L^2$ realization and
Casselman--Wallach extension from finite vectors, as well as the
archimedean up/down pairing and chain-level laws. These are separate
analytic assertions. The equivalence above neither assumes their
conclusions for the definite end nor gives an independent proof of
these two inner assertions.
% END E77A6
